\chapter{Sound and music}

Of everything left in the library, this is the one that changes how finished your
game feels for the least work. Three function calls and the mini Doom stops being
a tech demo.

\section{One line before anything}

\begin{lstlisting}
InitAudioDevice();      // right after InitWindow
...
CloseAudioDevice();     // right before CloseWindow
\end{lstlisting}

Nothing in raudio works before that call, exactly like \lstinline|InitWindow| for
graphics. If sounds silently do nothing, check this first.

\section{The four things raudio gives you}

\begin{center}
\begin{tabular}{@{}llp{6.4cm}@{}}
\toprule
\textbf{Type} & \textbf{Lives} & \textbf{Use it for} \\
\midrule
\lstinline|Wave| & in RAM, raw samples & building or editing audio yourself \\
\lstinline|Sound| & in RAM, decoded & short clips: gunshots, steps, clicks \\
\lstinline|Music| & streamed from disk & long tracks: the soundtrack \\
\lstinline|AudioStream| & you push samples & synthesisers, anything procedural \\
\bottomrule
\end{tabular}
\end{center}

The split between \lstinline|Sound| and \lstinline|Music| is about memory. A
three minute track decoded into RAM is around 30 MB; streamed, it costs a few
kilobytes of buffer. A footstep is 4 KB either way, so it may as well be resident
and instant.

\section{Sounds}

\begin{lstlisting}
Sound shot = LoadSound("shot.wav");      // once, outside the loop

if (IsMouseButtonPressed(MOUSE_BUTTON_LEFT)) PlaySound(shot);

UnloadSound(shot);                       // once, after the loop
\end{lstlisting}

Same load-outside, unload-after rule as textures. raylib reads \file{.wav},
\file{.ogg}, \file{.mp3}, \file{.flac}, \file{.qoa} and a few more.

\lstinline|PlaySound| does not wait and does not queue --- it fires and returns
immediately. Call it twice in one frame and you get two overlapping copies, which
is exactly what you want for gunfire.

\subsection{The trick that makes repeated sounds bearable}

A sound effect played identically fifty times in ten seconds becomes torture.
Real games vary it slightly on every shot:

\begin{lstlisting}
SetSoundPitch(shot, 0.85f + GetRandomValue(0, 30)/100.0f);
SetSoundPan(shot, GetRandomValue(20, 80)/100.0f);
PlaySound(shot);
\end{lstlisting}

Pitch is a multiplier, so 1.0 is unchanged and 0.5 is an octave down. Pan runs
0.0 (left) to 1.0 (right), with 0.5 centred.

\gotcha{\lstinline|SetSoundPitch| changes that \lstinline|Sound| object for
every future play, not just the next one. If you want two enemies to sound
different at the same time you need \lstinline|LoadSoundAlias()|, which shares
the samples but gives you a second set of playback settings, and costs almost
nothing.}

\section{Music}

\begin{lstlisting}
Music track = LoadMusicStream("theme.ogg");
track.looping = true;
PlayMusicStream(track);

while (!WindowShouldClose())
{
    UpdateMusicStream(track);     // <-- EVERY FRAME. Not optional.
    ...
}

UnloadMusicStream(track);
\end{lstlisting}

\gotcha{\lstinline|UpdateMusicStream()| is the single most forgotten call in
raylib. It refills the streaming buffer from disk. Miss it and your music plays
for a fraction of a second and stops --- no error, no warning, just silence. If
music ``does not work'', this is the reason nine times out of ten.}

Pause, resume and seek behave as you would expect, and two getters let you draw a
progress bar:

\begin{lstlisting}
PauseMusicStream(track);  ResumeMusicStream(track);  StopMusicStream(track);
SeekMusicStream(track, 30.0f);                        // jump to 30 seconds

float played = GetMusicTimePlayed(track);
float length = GetMusicTimeLength(track);
\end{lstlisting}

\section{Making sounds with no sound files}

You do not need an artist to get started. A \lstinline|Wave| is just numbers, and
a 16-bit mono wave is one \lstinline|short| per sample, from $-32768$ to $32767$:

\begin{lstlisting}
Wave wave = { 0 };
wave.sampleRate = 44100;      // samples per second
wave.sampleSize = 16;         // bits per sample
wave.channels   = 1;          // mono
wave.frameCount = 44100/5;    // a fifth of a second

short *samples = (short *)MemAlloc(wave.frameCount*sizeof(short));

float phase = 0.0f;
for (unsigned int i = 0; i < wave.frameCount; i++)
{
    float t = (float)i/wave.frameCount;        // 0..1 through the sound
    float hz = 1400.0f + (200.0f - 1400.0f)*t; // pitch slides DOWN
    phase += 2.0f*PI*hz/44100.0f;

    float envelope = (1.0f - t)*(1.0f - t);    // fade out, or it clicks
    samples[i] = (short)(sinf(phase)*envelope*22000.0f);
}

wave.data = samples;
Sound laser = LoadSoundFromWave(wave);
\end{lstlisting}

Three ideas are doing all the work there, and they are the whole of retro sound
design:

\begin{itemize}[leftmargin=*,itemsep=3pt]
  \item \textbf{A sine wave} gives you a pure tone. \lstinline|phase| accumulates
        and \lstinline|sinf(phase)| oscillates.
  \item \textbf{Sliding the pitch} turns a beep into a laser (down) or a power-up
        (up).
  \item \textbf{An envelope} --- multiplying by a value that falls to zero ---
        makes it a sound instead of a rude noise. Without a fade-out, the sound
        stops mid-oscillation and the speaker cone jumps, which you hear as a
        click.
\end{itemize}

Mix in random noise instead of a sine and you have an explosion. That is
literally what the \textbf{rFXGen} tool on raylib.com does, with sliders.

\gotcha{Use \lstinline|MemAlloc()| rather than \lstinline|malloc()| for wave
data. \lstinline|UnloadWave()| frees it with raylib's allocator, and mixing
allocators is the kind of bug that crashes only on some machines.}

\subsection{Wave and Sound are two separate allocations}

\lstinline|LoadSoundFromWave| copies the samples to the audio device. The
original \lstinline|Wave| is still yours, and still needs freeing:

\begin{lstlisting}
UnloadSound(laser);
UnloadWave(wave);      // both, not one
\end{lstlisting}

You can also write it out --- \lstinline|ExportWave(wave, "laser.wav")| --- which
is a genuinely useful way to build a sound in code, listen to it, and then ship
the file.

\section{Audio streams, for sound you generate live}

When the samples cannot exist ahead of time --- a synthesiser, an engine note
that follows the revs, voice chat --- you push them yourself:

\begin{lstlisting}
AudioStream stream = LoadAudioStream(44100, 16, 1);
PlayAudioStream(stream);

// in the loop:
if (IsAudioStreamProcessed(stream))     // the device wants more
{
    static short chunk[4096];
    for (int i = 0; i < 4096; i++) { phase += step; chunk[i] = ...; }
    UpdateAudioStream(stream, chunk, 4096);
}
\end{lstlisting}

The contract is: ask whether a buffer has been consumed, and if so, hand over the
next one. Miss the window and you hear a gap.

\section{Volume, and being polite}

\begin{lstlisting}
SetMasterVolume(0.6f);            // everything
SetSoundVolume(shot, 0.4f);       // one sound
SetMusicVolume(track, 0.5f);      // one track
\end{lstlisting}

Ship with the master volume somewhere around 0.5 and put it in an options menu.
Players will forgive almost anything except a game that is deafening on the first
launch.

\tryit{Module 9 synthesises a beep, a laser and an explosion in C, writes a
four second tune to a real \file{.wav} and streams it back, and lets you play a
live tone with the mouse controlling the pitch. No audio files
involved.}{./cheatsheet/build.sh 9}
